\section{Boundary comparison without boundary integrability}
\label{sec:boundary-comparison}

The minimum moment in \eqref{eq:minimum-moment} is the assumption A1 of
\cite{AhlbergHoffman}. A boundary expectation in the proof of their
Proposition~4.4 nevertheless requires care. Lemma~4.6 of that paper
constructs half-plane geodesics and then identifies the expectation of
an associated boundary Busemann increment. For a diagonal half-plane,
its boundary vertices have only two incident edges in the half-plane;
\eqref{eq:minimum-moment} need not give integrable restricted passage
times. We give a replacement for the expectation step in
\cite[Section~4.5]{AhlbergHoffman}. It proves the same uniqueness
conclusion using finite stationary corrections and bounded truncation.

For example, a continuous law with tail
\[
 \PP(t_e>t)=(1+t)^{-1/2}\bigl(\log(e+t)\bigr)^{-3/8},\qquad t\geq0,
\]
satisfies \eqref{eq:minimum-moment}, whereas the minimum of two
independent weights has infinite mean. Indeed, the relevant tail
integrals are comparable at infinity to
\(\int^\infty dt/(t(\log t)^{3/2})\) and
\(\int^\infty dt/(t(\log t)^{3/4})\), respectively. In
\(H=\{x_1+x_2\geq0\}\), the cost of leaving a boundary vertex is at
least that minimum. The proof below makes no assumption about the
integrability of the restricted Busemann increments.

We use the terminology of \cite[Definition~2.3]{AhlbergHoffman}:
a \emph{random coalescing geodesic} is a measurable choice of a ray
from the origin whose lattice translates coalesce almost surely.
Write \(\mathcal G(x)\) for its translated ray from \(x\), and
\(B_{\mathcal G}\) for their common full-plane Busemann function.
The full-plane linearity of this function is
\cite[Proposition~4.2]{AhlbergHoffman}, which precedes the boundary
argument. Let \(H\) be any of the eight axial or diagonal half-planes
used there, let \(z\) be a primitive lattice vector parallel to its
boundary, and let \(T_H\) be the passage metric restricted to \(H\).
Assume that \(\mathcal G\) eventually enters every translate of \(H\).
The finite-path argument in \Cref{sec:preliminaries} applies also to
the restricted lattice: its simple paths form a subcollection of the
full-plane paths, and the same nonatomic tie argument gives uniqueness.
Thus finite restricted minimizers exist and are unique, simultaneously
for these half-planes and their lattice translates. Positivity also
rules out loops in a minimizing walk. These are pathwise assertions,
not integrability assertions about \(T_H\).

\begin{lemma}[Finite boundary correction]\label{lem:boundary-correction}
For \(x\in H\cap\Z^2\), the limit
\[
 D_{\mathcal G}(x)
 =\lim_{j\to\infty}\bigl(T_H(x,v_j)-T(x,v_j)\bigr)
\]
exists and is finite and nonnegative, where \((v_j)\) is any ray of
\(\mathcal G\) and the targets are sufficiently far along its
\(H\)-contained tail. The limit is independent of the target ray.
The variables \(D_{\mathcal G}(kz)\), \(k\in\Z\), are stationary
under boundary translations, and the restricted Busemann limit satisfies
\begin{equation}\label{eq:boundary-correction}
 B_{\mathcal G}^H(x,y)
 =B_{\mathcal G}(x,y)+D_{\mathcal G}(x)-D_{\mathcal G}(y).
\end{equation}
\end{lemma}
\begin{proof}
Choose a vertex \(a\) after the last exit of the target ray from \(H\).
The segment from \(a\) to each later \(v_j\) minimizes for both metrics,
so \(T_H(a,v_j)=T(a,v_j)\). Each of the sequences
\[
 T_H(x,v_j)-T(a,v_j),\qquad T(x,v_j)-T(a,v_j)
\]
is decreasing by the triangle inequality along this calibrated tail.
They are bounded in absolute value by \(T_H(x,a)\) and \(T(x,a)\),
respectively. These bounds are finite because the half-plane lattice
is connected and all edge weights are finite. Their limits therefore
exist and are finite. Their difference is nonnegative, since
\(T_H\geq T\), and defines \(D_{\mathcal G}(x)\).

Any two rays of the family have a common tail, so changing the target
ray leaves the limit unchanged. A translation by \(kz\) preserves
\(H\) and replaces the target ray by another translated ray of the
same family. This proves covariance, hence stationarity. No choice
of the auxiliary anchor \(a\) enters the definition of \(D_{\mathcal G}\).
Measurability follows from the measurable ray choice and the countable
limits defining the distances. Subtracting the two distance limits
at \(x\) and \(y\) proves \eqref{eq:boundary-correction} and the
existence of \(B_{\mathcal G}^H\).
\end{proof}

\begin{lemma}[A nonnegative stationary difference]\label{lem:zero-boundary-difference}
Let \(\Delta(k,l)\), \(k,l\in\Z\), be a finite, additive stationary
process, with \(\Delta(k,k+1)\geq0\) almost surely. If
\(\Delta(0,n)/n\to0\) in probability, then every increment is zero
almost surely.
\end{lemma}
\begin{proof}
Set \(X_k=\min\{\Delta(k,k+1),1\}\). Then
\[
 0\leq A_n:=\frac1n\sum_{k=0}^{n-1}X_k
 \leq\frac{\Delta(0,n)}n,\qquad A_n\leq1.
\]
Thus \(A_n\to0\) in probability and
\(\EE A_n\leq\eps+\PP(A_n>\eps)\to\eps\) for every \(\eps>0\).
Stationarity gives \(\EE A_n=\EE X_0\) for all \(n\), so
\(X_0=0\) almost surely. Stationarity, countability and additivity
prove the assertion for every pair. Neither integrability of
\(\Delta\) nor ergodicity is needed.
\end{proof}

\begin{proposition}[The fixed-functional boundary comparison]
\label{prop:boundary-uniqueness}
Two random coalescing geodesics with the same full-plane Busemann
linear functional agree almost surely under the assumptions of
\Cref{thm:main}.
\end{proposition}
\begin{proof}
Let the two families be \(\mathcal G,\mathcal G'\), with common
functional \(\rho\). By
\cite[Proposition~4.3]{AhlbergHoffman}, their contact sector is
contained in the interior of one of the eight half-planes \(H\).
Both families eventually enter every translate of \(H\).

We first recall the geometric part of
\cite[Lemma~4.6, Section~4.4]{AhlbergHoffman}, which does not use its
subsequent expectation assertion. There is positive probability that
the ray from the origin stays in \(H\): on an eventual-entrance ray,
a vertex minimizing the normal coordinate starts a tail contained in
its translated half-plane. Finite uniqueness and coalescence identify
this suffix with the selected ray from that vertex. Countability of the
vertices and covariance give positive probability at a fixed vertex.
Boundary translations of an iid environment are ergodic, also for
diagonal boundaries. Hence there are such starting vertices in both
directions along the boundary.

Choose two of these roots on opposite sides of the origin. The two
contained rays coalesce, and their portions up to their first common
vertex form a finite simple connecting arc. If the origin lies on this
arc, its selected suffix is already contained in \(H\). Otherwise,
choose consecutive boundary contacts of the arc on opposite sides of
the origin. The intervening subarc, together with the boundary segment,
encloses a bounded region adjacent to the origin. Choose a vertex
\(w\) on the common tail of the two bounding rays and the target ray,
outside this region and beyond the connecting arc. Every restricted
minimizer from the origin to a target beyond \(w\) must meet the
separating subarc. From that meeting point, the bounding-ray suffix to
the target lies in \(H\) and is a full-plane minimizer. Restricted
uniqueness forces this suffix, hence passage through \(w\). The prefix
to \(w\) is then the fixed restricted minimizer to \(w\). This proves
stabilization and coalescence with the full-plane family, including when
the bounding rays revisit the boundary. Applying the construction at
every boundary root gives \(\mathcal G_H,\mathcal G'_H\); their common
Busemann functions on the boundary are exactly the restricted limits
in \Cref{lem:boundary-correction}.

Suppose that the families differ with positive probability. Because
each restricted family coalesces, the event that the two families
coalesce is invariant under boundary translations. If they coalesce,
finite uniqueness makes their rays from any common root identical.
Boundary ergodicity therefore makes distinctness an almost-sure event.
The same invariance and ergodicity make their asymptotic planar order
deterministic, as in \cite[Section~4.5]{AhlbergHoffman}. Orient the boundary
vector \(z\) so that the finite crossing comparison in equation~(12)
of that source gives
\begin{equation}\label{eq:boundary-nonnegative}
 \Delta(k,l):=B_{\mathcal G}^H(kz,lz)
             -B_{\mathcal G'}^H(kz,lz)\geq0,
 \qquad k<l.
\end{equation}
For clarity, the comparison chooses a crossing of
\(\mathcal G_H(kz)\) with \(\mathcal G'_H(lz)\) and common-tail
vertices \(x\) of the primed rays and \(y\) of the unprimed rays,
both beyond the crossing.
Swapping the two pieces at the crossing gives walks \(\pi_x\) from
\(kz\) to \(x\) and \(\pi_y\) from \(lz\) to \(y\), and
\[
 \Delta(k,l)
 =\bigl(T_H(\pi_y)-T_H(lz,y)\bigr)
  +\bigl(T_H(\pi_x)-T_H(kz,x)\bigr).
\]
Both summands are nonnegative finite excess costs. This identity is
purely geometric and does not use expectations.

Put \(Z(x)=D_{\mathcal G}(x)-D_{\mathcal G'}(x)\).
Equation~\eqref{eq:boundary-correction} gives
\[
 \Delta(0,n)=B_{\mathcal G}(0,nz)-B_{\mathcal G'}(0,nz)
                    +Z(0)-Z(nz).
\]
The first difference divided by \(n\) tends to zero almost surely,
by the full-plane linearity and the equality of the two functionals.
For each \(\eps>0\), finiteness and stationarity give
\[
 \PP(|Z(nz)|>n\eps)=\PP(|Z(0)|>n\eps)\longrightarrow0.
\]
Consequently \(\Delta(0,n)/n\to0\) in probability. The process
\(\Delta\) is additive and stationary under boundary shifts, so
\Cref{lem:zero-boundary-difference} and
\eqref{eq:boundary-nonnegative} imply \(\Delta(k,l)=0\) for every
pair, almost surely.

The final uniqueness argument is the finite one in
\cite[Section~4.5]{AhlbergHoffman}. Distinct same-root restricted
rays have a finite common initial portion. Choose a deterministic
bound on its length that holds with positive probability; boundary
ergodicity gives two such roots farther apart than twice this bound.
In the crossing comparison just described, at least one swapped walk
then differs from the corresponding minimizing path. Indeed, if both
swapped walks were their unique geodesics, their crossing vertex would
lie in the common initial portion from each root. If these portions
have at most \(m\) edges, the lattice distance between the roots would
be at most \(2m\), contrary to their choice. Since both nonnegative
excesses sum to zero, a distinct swapped walk is also minimizing,
contrary to uniqueness. Thus the restricted families
agree. Their coalescence with the full-plane families and finite
uniqueness imply \(\mathcal G=\mathcal G'\).
\end{proof}

This supplies the conclusion of
\cite[Proposition~4.4]{AhlbergHoffman} without the boundary-expectation
step. The other later direct use of Lemma~4.6, in the proof of
Lemma~9.6 of that source, uses only the geometric construction and
coalescence. Thus the structural input in \Cref{thm:structure} retains
its stated minimum-moment scope; no integrability of a diagonal
boundary Busemann function is inserted.
